\section{A Qualified Test on a Language Model}\label{sec:real}

\subsection{Setup}
The primary model is an 848K-parameter, 4-layer, 128-dimensional character-level Transformer with a 256-character context, trained on Austen and Shelley (checkpoint \textsc{cosine}; Appendix~\ref{sec:protocol}). Every experiment uses the 15 \textsc{test} prompts (Appendix~\ref{app:prompts}), 10 seeds per prompt, generations of exactly 500 characters, and temperature $T=0.10$, a regime in which this model loops: in a 1000-character temperature sweep the loop rate falls from 96\% at $T=0.05$ to 0\% by $T=0.30$ (Appendix~\ref{sec:protocol}). The reference $P$ is plain temperature sampling at $T=0.10$. Recurrence risk uses $n$-gram orders 3--6. The loop event is the earliest position at which a span of length $P\in[2,60]$ repeats five times consecutively (a separate threshold of 8 for single-character runs). It is validated on seven held-out corpora with a pooled false-positive rate of 1.3\% (95\% CI [0.9\%, 1.9\%]), concentrated in one corpus (\emph{Dracula}, 8.7\%, epistolary formulas), and it is purely syntactic (Appendix~\ref{sec:event}); loop counts are counts of generations in which the detector fires, with that caveat attached. Decoders: the three recurrence-risk modes (Section~\ref{sec:setup}), repetition penalty (1.3), hard no-repeat 4-gram blocking, SuffixMatch ($\alpha=3$), and an FSD-style decoder and an LZ Penalty, both adapted from the published methods (Appendix~\ref{app:baselines}). Hard no-repeat blocking forbids the measured event by construction, so its zero loops are not a finding and it is excluded from every ranking and test (Appendix~\ref{ssec:hard-repeat}).

\subsection{Observed loop counts and per-step distortion}\label{ssec:common-kl}
{{TABLE:primary}}
Table~\ref{tab:primary} scores every decoder against the same reference and the same risk definition, by exposing each decoder's per-step distribution. The reference loops in 43 of 150 generations. Observed counts: \textsc{rr-dual} 0/150 at both $\varepsilon$ (95\% upper limit 2.4\%), \textsc{rr-adaptive} 0/150, LZ Penalty 0/150, SuffixMatch 0/150, FSD-style 3/150, \textsc{rr-fixed} 4/150, repetition penalty 7/150. These are counts at one temperature on one model, not guarantees. Per-step distortion differs by a factor of more than 100 among the decoders that reach zero observed loops (0.096 bits for \textsc{rr-adaptive} to 12.7 for \textsc{rr-dual}). The per-step guarantee of the dual mode held in every step (feasibility 100.0\%, zero maximum violation at the stored precision; Appendix~\ref{app:baselines}), so in this experiment the fixed-budget decoder is never infeasible, the situation of the $R=3$ system of Section~\ref{ssec:infeas}. RMST and paired comparisons are in Appendix~\ref{sec:protocol}: the exact test does not separate \textsc{rr-fixed} from the decoders with zero observed loops ($p=0.125$ at 15 prompts).

\subsection{A fixed price against a fixed budget}\label{ssec:frontier}
The theory suggests that a fixed-price rule should reach a given reduction of the loop probability at lower KL than a fixed-budget rule, and that the budget rule realizes widely varying prices. In the recurrence-risk decoder the two are the \textsc{fixed} mode (one $\lambda$ for the whole generation) and the \textsc{dual} mode (one $\varepsilon$ re-solved at every step). We sweep both on the same prompts, seeds, risk definition and detector, with $\lambda\in\{0.25,0.5,1,2,3,4,6,8,12\}$ and $\varepsilon\in\{1,0.95,0.9,0.85,0.8,0.75,0.7,0.65,0.6,0.5,0.3,0.2,0.1,0.05,0.01\}$, and record for every generation its loop event, its path KL and the per-step multiplier trace of the decoder. The budget $\varepsilon=1$ is a control: the constraint $\E_q[\mathrm{risk}]\le1$ is never active, so the dual mode must reproduce the reference decoder, which we check generation by generation. The comparison is descriptive: each setting has 75 generations, so an observed loop count does not establish equality of true loop probabilities, and we compare cost at similar observed counts with the intervals stated.

{{FRONTIER_RESULTS}}

\subsection{How much of the path KL moves the loop probability}\label{ssec:eventlevel}
Section~\ref{ssec:bernoulli-bound} splits a decoder's path divergence into the term $d_{\mathrm{Ber}}(Q(E)\Vert P(E))$ needed to move the loop probability and a remainder, the distortion of the conditional shape of the generation. We estimate both for every decoder from the sample-level records of the run behind Table~\ref{tab:primary} (for each generation: its prompt, loop event and total path KL).

\textbf{The comparison is conditioned on the prompt.} The decomposition applies prompt by prompt, $\KL(Q_x\Vert P_x)\ge d_{\mathrm{Ber}}(Q_x(E)\Vert P_x(E))$, and the quantities reported are averages over the 15 prompts with weights $1/15$. Pooling the prompts first would use the pooled loop rates and, by joint convexity of $d_{\mathrm{Ber}}$, can only lower the bound, overstating the distortion that is not needed; the reference loop count differs by prompt (1 to 6 of 10 generations), but the difference is small here (at a decoder loop fraction of zero the plug-in event term is 0.528 bits prompt by prompt and 0.487 bits pooled). $P_x(E)$ and $Q_x(E)$ are loop fractions over the 10 seeds of one prompt; the plug-in $d_{\mathrm{Ber}}$ is infinite whenever a fraction is 0 or 1 and the other is not, so the primary estimate smooths both fractions as $(k+\tfrac12)/(n+1)$, with plug-in and Laplace ($(k+1)/(n+2)$) smoothings as a sensitivity range. Intervals are percentile intervals from resampling whole prompts (10{,}000 resamples), which keeps each generation's loop event and path KL together. Everything is a sample estimate from 150 generations per decoder; it is descriptive and has no guaranteed coverage.

{{TABLE:pcb}}
For every decoder the estimated event term lies between 0.20 and 0.53 bits under all three smoothings, whereas the measured path KL ranges from 19 to over 6000 bits (Table~\ref{tab:pcb}). At least 98\% of what each decoder spends is estimated to be distortion of the conditional shape of the generation and not the cost of moving the loop probability: 98.4\% for \textsc{rr-fixed}, above 99.9\% for FSD-style, no-repeat blocking and \textsc{rr-dual}. The event term itself is uncertain by about $\pm0.15$ bits because it rests on 10 generations per prompt, which does not affect this reading.

\textbf{What this does and does not say.} The excess is measured against the cost of moving the loop probability by reweighting whole sequences by whether they loop. That bound is attained only by a rule that has the hazard at every step, or by rejecting whole generations; Section~\ref{sec:small} shows that even with the exact hazard a fixed budget spends multiples of it. The excess of a decoder is therefore not a saving available to that decoder, and we do not claim that 98\% of any decoder's distortion could be removed. The estimate rules out one reading, that the large KL of the heavy decoders is spent on lowering the loop probability, and it says nothing about text quality.

\subsection{What transfers and what remains a hypothesis}\label{ssec:transfer}
{{TABLE:transfer}}
The prompt-conditional analysis was run on the primary model only, the only model for which sample-level path KL is stored. The other models (a character model trained on non-fiction, GPT-2, OLMo-2-0425-1B, Mistral-7B-v0.3; Appendix~\ref{sec:generalization}) are compared at a compact scale and show that the heuristics' effectiveness varies from model to model: \textsc{rr-fixed} has 4/150 loops on the primary model and 77/150 on the non-fiction model, and the LZ Penalty and the FSD-style decoder reach zero observed loops on some models and not on others. The dual mode's per-step feasibility was exactly 100.0\% on every model. Qwen2.5-3B showed no repetition regime at the budget tested and is not compared.
